\documentclass[11pt,a4paper]{article}

\usepackage[utf8]{inputenc}
\usepackage[T1]{fontenc}
\usepackage[polish]{babel}
\usepackage[margin=2.5cm]{geometry}
\usepackage{amsmath,amssymb}
\usepackage{graphicx}
\usepackage{booktabs}
\usepackage{longtable}
\usepackage{array}
\usepackage{xcolor}
\usepackage{tikz}
\usepackage{hyperref}

\usetikzlibrary{positioning,decorations.pathreplacing}

\hypersetup{colorlinks=true, linkcolor=blue, urlcolor=blue, citecolor=blue}

% spis treści tylko do poziomu podsekcji
\setcounter{tocdepth}{2}

% notatka robocza — dane w trakcie zbierania
\newcommand{\wtoku}[1]{\textcolor{red}{\textbf{[W TOKU: #1]}}}

\title{%
    Topology-aware ewaluacja mitygacji błędów\\
    i narzutu kubitów logicznych na sprzęcie nadprzewodzącym\\
    \large Sprawozdanie z Miesiąca 1
}
\author{Karina Leśkiewicz}
\date{Sierpień 2026}

\begin{document}

\maketitle
\thispagestyle{empty}
\clearpage

% =====================================================================
\section*{Streszczenie}
\addcontentsline{toc}{section}{Streszczenie}
% =====================================================================

Badam, w którym momencie teoretyczny zysk z mitygacji błędów zostaje zniwelowany przez
fizyczne ograniczenia sprzętu --- w szczególności przez narzut routingu (bramki SWAP)
wymuszony topologią procesora. Jako zmienną niezależną przyjęłam liczbę bramek
dwukubitowych po kompilacji (transpiled CZ), a jako sygnał jakości --- błąd wartości
oczekiwanej $|\langle O\rangle_{\text{noisy}} - \langle O\rangle_{\text{ideal}}|$,
dokładność klasyfikatora oraz fidelity obwodu.

W Miesiącu 1 zbudowałam powtarzalny zestaw narzędzi: trening w kalibrowanym symulatorze
szumu, a następnie inferencję z REM${}+{}$ZNE na sprzęcie ODRA (IQM Spark, topologia
gwiazdy). Uzyskałam kompletną routing ladder dla zbioru \texttt{two\_curves} na
rzeczywistym QPU --- zarówno błąd wartości oczekiwanej, jak i mirror fidelity --- oraz
drugą, niezależną routing ladder (ansätze Banknote) o porównywalnej dokładności
bezszumowej i trzykrotnym rozrzucie kosztu CZ. Zaobserwowałam, że skuteczność ekstrapolacji
ZNE maleje monotonicznie z kosztem CZ i staje się znikoma w okolicach $85$--$90$ bramek CZ,
przy jednoczesnym odzysku dokładności klasyfikatora na sprzęcie tam, gdzie szum zdążył
przerzucić znak wartości oczekiwanej. Wykazałam również, że Dynamical Decoupling
w natywnej konfiguracji urządzenia \emph{pogarsza} wszystkie mierzone wielkości, i opisałam
to jako wynik negatywny. Pełne kodowanie logiczne (QEC) odrzuciłam z uzasadnionych powodów
sprzętowych; stosuję mitygację błędów (QEM).

\clearpage
\tableofcontents
\clearpage

% =====================================================================
\section{Wprowadzenie i motywacja}
% =====================================================================

\paragraph{Dlaczego mitygacja błędów (QEM).} Obecne procesory kwantowe działają
w reżimie NISQ (Noisy Intermediate-Scale Quantum): mają dziesiątki do setek kubitów, ale
błąd pojedynczej bramki dwukubitowej (rzędu $1$--$2\%$) i ograniczone czasy koherencji
uniemożliwiają wykonywanie głębokich obwodów bez akumulacji szumu. Docelowym rozwiązaniem
jest korekcja błędów (QEC), która koduje jeden kubit logiczny w wiele fizycznych i pozwala
na dowolnie długie obliczenia --- lecz wymaga błędu poniżej progu oraz narzutu zasobów o
rzędy wielkości, czego dzisiejszy sprzęt nie spełnia (szczegółowo w sekcji~\ref{sec:qec}).
Powstaje luka: chcemy uruchamiać użyteczne algorytmy \emph{już teraz}, na sprzęcie, którego
nie stać na pełną korekcję.

Mitygacja błędów (Quantum Error Mitigation, QEM) wypełnia tę lukę. Zamiast korygować błędy
w trakcie obliczeń, QEM \emph{redukuje ich wpływ na końcową wartość oczekiwaną} metodami
statystycznymi i post-processingiem (Zero-Noise Extrapolation, Readout Error Mitigation,
Dynamical Decoupling), bez dodatkowych kubitów i bez feed-forward. Kosztem jest zwiększona
wariancja, czyli więcej powtórzeń pomiaru. Dlatego QEM jest dziś głównym praktycznym narzędziem otrzymywania wiarygodnych
wyników z maszyn NISQ i przedmiotem intensywnych badań~\cite{cai2023}: to pomost między erą
szumu a erą korekcji. Kluczowe pytanie, które QEM stawia, jest jednak empiryczne --- \emph{do
jakiego poziomu szumu (kosztu obwodu) mitygacja realnie pomaga, zanim jej narzut wariancji
przeważy zysk}. To pytanie jest osią niniejszego projektu.

\paragraph{Rola topologii sprzętu.}
Odwzorowanie logicznych kubitów algorytmu na fizyczne kubity procesora jest ograniczone
jego topologią. Gdy algorytm wymaga oddziaływania między kubitami, które nie są fizycznie
połączone, kompilator wstawia operacje SWAP. Pojedynczy SWAP rozkłada się na trzy bramki
CNOT (lub natywny odpowiednik dwukubitowy), co zwiększa głębokość obwodu, akumuluje błąd
dwukubitowy i wydłuża okna bezczynności podatne na dekoherencję.

\paragraph{Cel główny.} Empiryczne ustalenie, przy jakim koszcie dwukubitowym (po
kompilacji) akumulacja szumu z routingu wypycha urządzenie poza reżim, w którym
ekstrapolacja ZNE pozostaje wiarygodna --- czyli gdzie mitygacja przestaje pomagać.

\paragraph{Pytania badawcze.}
\begin{enumerate}
\item Jaka jest bazowa jakość wybranego obwodu na kubitach fizycznych (wartość oczekiwana,
      fidelity)?
\item Czy wariant logiczny/zakodowany poprawia wynik pomiarowy względem fizycznego?
\item Jak mapowanie optymalne i celowo niekorzystne zmienia głębokość, liczbę SWAP
      i jakość wyjścia?
\item Czy istnieje mierzalna korelacja między fizyczną odległością kubitów, liczbą SWAP
      i degradacją jakości modelu?
\end{enumerate}

\paragraph{Sprzęt (trójpoziomowy dostęp).}

\begin{center}
\begin{tabular}{@{}llccl@{}}
\toprule
\textbf{Poziom} & \textbf{Urządzenie} & \textbf{Kubity} & \textbf{Topologia} & \textbf{Rola} \\
\midrule
A & ODRA (IQM Spark) & 5  & gwiazda        & podstawowy; generator SWAP \\
B & IQM Garnet       & 20 & krata kwadratowa & walidacja pośrednia \\
C & IQM Emerald      & 54 & krata kwadratowa & walidacja wielkoskalowa \\
\bottomrule
\end{tabular}
\end{center}

\begin{figure}[htbp]
\centering
% --- (a) ODRA: gwiazda 5 kubitów ---
\begin{minipage}[b]{0.30\textwidth}
\centering
\begin{tikzpicture}[scale=0.95,
    qb/.style={circle,draw,minimum size=6.5mm,inner sep=0pt,font=\tiny}]
  \node[qb,fill=black!10] (hub) at (0,0) {QB3};
  \node[qb] (q1) at (-1.1, 1.0) {QB1};
  \node[qb] (q2) at ( 1.1, 1.0) {QB2};
  \node[qb] (q4) at (-1.1,-1.0) {QB4};
  \node[qb] (q5) at ( 1.1,-1.0) {QB5};
  \draw (hub)--(q1); \draw (hub)--(q2); \draw (hub)--(q4); \draw (hub)--(q5);
\end{tikzpicture}
\\[2mm]{\small (a) ODRA (IQM Spark), 5 kubitów}
\end{minipage}
\hfill
% --- (b) Garnet: 20 kubitów, krata 4x5 ---
\begin{minipage}[b]{0.30\textwidth}
\centering
\begin{tikzpicture}[scale=0.42,
    qb/.style={circle,draw,minimum size=3.6mm,inner sep=0pt}]
  \foreach \x in {0,...,3}{\foreach \y in {0,...,4}{
      \node[qb] (g\x\y) at (\x*1.5,\y*1.5) {};}}
  \foreach \x in {0,...,2}{\foreach \y in {0,...,4}{
      \draw (g\x\y)--(g\the\numexpr\x+1\relax\y);}}
  \foreach \x in {0,...,3}{\foreach \y in {0,...,3}{
      \draw (g\x\y)--(g\x\the\numexpr\y+1\relax);}}
\end{tikzpicture}
\\[2mm]{\small (b) IQM Garnet, 20 kubitów}
\end{minipage}
\hfill
% --- (c) Emerald: 54 kubity, krata 6x9 ---
\begin{minipage}[b]{0.34\textwidth}
\centering
\begin{tikzpicture}[scale=0.34,
    qb/.style={circle,draw,minimum size=3mm,inner sep=0pt}]
  \foreach \x in {0,...,5}{\foreach \y in {0,...,8}{
      \node[qb] (e\x\y) at (\x*1.5,\y*1.5) {};}}
  \foreach \x in {0,...,4}{\foreach \y in {0,...,8}{
      \draw (e\x\y)--(e\the\numexpr\x+1\relax\y);}}
  \foreach \x in {0,...,5}{\foreach \y in {0,...,7}{
      \draw (e\x\y)--(e\x\the\numexpr\y+1\relax);}}
\end{tikzpicture}
\\[2mm]{\small (c) IQM Emerald, 54 kubity}
\end{minipage}
\caption{Topologie trzech urządzeń z planu badawczego. (a) ODRA --- gwiazda z hubem QB3:
każda interakcja między kubitami zewnętrznymi musi przejść przez hub, generując SWAP.
(b) Garnet i (c) Emerald --- kraty kwadratowe o stopniu wierzchołka do 4, w których
sąsiedztwo jest znacznie bogatsze, więc ten sam obwód logiczny wymaga mniej operacji SWAP.
Panele (b) i (c) są schematami poglądowymi o właściwej liczbie kubitów; służą porównaniu
klasy łączności, a nie odwzorowaniu dokładnego układu sprzężeń.}
\end{figure}

% =====================================================================
\section{Podstawy teoretyczne}
% =====================================================================

W tej sekcji przedstawiam metody, które badałam:
Dynamical Decoupling (DD), Readout Error Mitigation (REM) oraz Zero-Noise Extrapolation
(ZNE), wraz z ramami bias/wariancja, w których działa mitygacja błędów~\cite{cai2023}.

\subsection{Ramy: bias, wariancja i sens mitygacji}

Skupiam się na średniokwadratowym odchyleniu estymatora $\hat{O}$ od wartości prawdziwej
$\mathrm{Tr}[O\rho_0]$:
\begin{equation}
\mathrm{MSE}[\hat{O}] = \mathbb{E}\!\left[(\hat{O}-\mathrm{Tr}[O\rho_0])^2\right]
= \mathrm{Bias}[\hat{O}]^2 + \mathrm{Var}[\hat{O}].
\end{equation}
Przy dużej liczbie wykonań MSE jest ograniczone przez bias --- błąd systematyczny, którego
nie usuwa zwiększanie liczby powtórzeń. Mitygacja konstruuje estymator $\hat{O}_{em}$
o mniejszym biasie, $|\mathrm{Bias}[\hat{O}_{em}]| \le |\mathrm{Bias}[\hat{O}_\rho]|$, lecz
\begin{equation}
C_{em} = \frac{\mathrm{Var}[\hat{O}_{em}]}{\mathrm{Var}[\hat{O}_\rho]}.
\end{equation}

To wyjaśnia, dlaczego ciągły błąd wartości oczekiwanej jest czulszym miernikiem niż
accuracy: mitygacja redukuje bias $\langle O\rangle$, ale przy małej
próbie blisko granicy decyzyjnej dodana wariancja może nieznacznie przesuwać
dokładność~\cite{cai2023}. Ta zależność jest bezpośrednio widoczna w wynikach
(sekcja~\ref{sec:tc-acc}): ZNE redukuje błąd $\langle O\rangle$ nawet 4,6-krotnie, nie
zmieniając przy tym dokładności.

\subsection{Dynamical Decoupling (DD)}

DD to technika sterowania w otwartej pętli, tłumiąca dekoherencję kubitów bezczynnych
($T_1$, $T_2$) wynikającą z oddziaływania układ--otoczenie. W kontekście kompilacji
topology-aware bramki SWAP znacząco zwiększają głębokość, pozostawiając nieuczestniczące
kubity bezczynne przez długi czas.

Dla Hamiltonianu układu i otoczenia $H = H_S\otimes I_E + I_S\otimes H_E + H_{int}$
wprowadzamy sterujący $H_c(t)$ jako ciąg natychmiastowych, równoodległych impulsów. Dla $N$
impulsów $\{P_k\}$ w odstępach $\tau$ operator ewolucji to
\begin{equation}
U(T) = \prod_{k=1}^{N}\left(P_k\, e^{-iH\tau}\right),
\end{equation}
a przy odpowiednim doborze sekwencji średni efekt oddziaływania znika w pierwszym rzędzie
rozwinięcia Magnusa:
\begin{equation}
H^{(0)}_{int} = \frac{1}{T}\sum_{k=1}^{N} P_k^{\dagger} H_{int} P_k \approx 0.
\end{equation}

\paragraph{Automatyczne wstawianie DD przez kompilator.}
(1) transformacja obwodu do DAG i mapowanie na topologię, czyli wstawianie SWAP;
(2) pobranie kalibracji: czasy bramek w ns i rozdzielczość zegara $dt$;
(3) harmonogram ALAP/ASAP z przypisaniem $t_{start}$, $t_{end}$, przy czym ALAP minimalizuje
czas w superpozycji;
(4) ekstrakcja okien bezczynności $T_{idle}$;
(5) padding: od $T_{idle}$ odejmujemy czas impulsów, a resztę dzielimy na równe $\tau$ zgrane
z zegarem;
(6) wstrzyknięcie impulsów.
Na sprzęcie IQM DD jest natywną opcją kompilacji i korzystam z wbudowanej strategii
urządzenia; w symulatorze modeluję je zamianą tempa dekoherencji
$T_2^{*} \rightarrow T_2^{\text{echo}}$ (echo Hahna).

\subsection{Readout Error Mitigation (REM)}

REM to klasyczne przetwarzanie po wykonaniu obwodu. W kubitach nadprzewodzących fidelity
pomiaru jest ograniczona, więc sprzęt myli stan końcowy. Błędy odczytu działają
jak klasyczny proces Markowa na rozkładzie prawdopodobieństw. Konstruujemy macierz
przypisania $A$ o wymiarze $2^n\times 2^n$:
\begin{equation}
A_{i|j} = P(\text{zmierzono } i \mid \text{przygotowano idealnie } j),
\qquad \vec{p}_{\text{noisy}} = A\,\vec{p}_{\text{ideal}}.
\end{equation}
Zwykłe odwrócenie $\vec{p}_{\text{mit}} = A^{-1}\vec{p}_{\text{noisy}}$ przy skończonej
liczbie shotów daje wyniki niefizyczne (ujemne prawdopodobieństwa), dlatego stosuje się
ograniczony problem najmniejszych kwadratów:
\begin{equation}
\min_{\vec{x}} \|A\vec{x} - \vec{p}_{\text{noisy}}\|_2
\quad \text{przy} \quad x_i \ge 0,\ \textstyle\sum_i x_i = 1.
\end{equation}
W mojej implementacji (obserwable $Z/I$) używam dokładnego, per-kubitowego modelu
symetrycznego bit-flip: symetryczny błąd $p$ tłumi $\langle Z\rangle$ o czynnik $(1-2p)$,
więc nieobciążony estymator to
$\langle Z\rangle_{\text{mit}} = \langle Z\rangle_{\text{meas}}/(1-2p)$.
Podanie wolnych-od-błędu-odczytu wartości oczekiwanych do ZNE jest istotne: pozwala ZNE
skupić się wyłącznie na szumie bramek, bez klasowego biasu pomiaru~\cite{koh2024,cai2023}.

\subsection{Zero-Noise Extrapolation (ZNE)}

ZNE wykorzystuje stany zaszumione przy różnych poziomach błędu $\lambda$: stan $\rho_\lambda$,
a $\mathrm{Tr}[O\rho_\lambda]$ traktujemy jako funkcję $\lambda$. Sondujemy przy zwiększonych
$\{\lambda_m\}$ i dażymy do $\lambda\to 0$. Poziom szumu zwiększamy cyfrowo przez
unitary folding $U \to U(U^{\dagger}U)^k$ dla $\lambda = 2k+1 \in \{1,3,5\}$.

\paragraph{Ekstrapolacja Richardsona.} Zakładamy, że wartość oczekiwana jest funkcją
wielomianową poziomu szumu, $\mathrm{Tr}[O\rho_\lambda] = a_0 + a_1\lambda + \dots +
a_{M-1}\lambda^{M-1}$. Interesuje nas wyraz wolny $a_0$, czyli wartość w $\lambda = 0$. Mając
pomiary w $M$ punktach $\lambda_1,\dots,\lambda_M$ przeprowadzamy przez nie \emph{dokładnie
jeden} wielomian stopnia $M-1$ i odczytujemy jego wartość w zerze. Jest to interpolacja
Lagrange'a wyliczona w punkcie $0$, co daje kombinację liniową pomiarów o stałych
współczynnikach:
\begin{equation}
\mathbb{E}[\hat{O}_{em}] = \sum_{m=1}^{M} \mathrm{Tr}[O\rho_{\lambda_m}]
\prod_{k\neq m}\frac{\lambda_k}{\lambda_k-\lambda_m},
\qquad
C_{em} \sim \left(\sum_m \left|\prod_{k\neq m}\frac{\lambda_k}{\lambda_k-\lambda_m}\right|\right)^{2}.
\end{equation}
Istotą metody jest to, że współczynniki są \emph{naprzemienne co do znaku} i większe od
jedności, więc ekstrapolacja jest różnicą dużych liczb. Stąd bierze się narzut wariancji:
$C_{em}$ rośnie, gdy $\lambda_m$ jest za duże albo gdy odstępy $|\lambda_m-\lambda_k|$ są za
małe. Gdy $C_{em}$ eksploduje, ekstrapolacja staje się niewykonalna --- i to jest sam w sobie
wskaźnik załamania metody. Dla dużych $\lambda$ stosuję wariant wykładniczy. Ideę
„sztucznego zwiększania i ekstrapolacji błędu” wprowadzili Li i Benjamin~\cite{li2017};
porównuję warianty liniowy, Richardsona i wykładniczy.

% =====================================================================
\section{Metodologia i uzasadnienia}
% =====================================================================

\subsection{VQC, obserwable i metryki jakości}

Model realizuje odwzorowanie parametryzowanym obwodem $U(x,\theta)$; predykcja to znak
wartości oczekiwanej obserwabli:
\begin{equation}
f(x,\theta) = \langle O\rangle = \langle 0|U^{\dagger} O\, U|0\rangle,
\qquad \hat{y} = \mathrm{sgn}\,\langle O\rangle,
\end{equation}
z obserwablą mean value $O_{\text{mean}} = \frac{1}{n}\sum_q Z_q$ (\texttt{two\_curves}) lub
pojedynczym $Z$ na kubicie readout (Banknote).

\paragraph{Błąd wartości oczekiwanej.} Definiuję go jako
$\varepsilon = |\langle O\rangle_{\text{noisy}} - \langle O\rangle_{\text{ideal}}|$, gdzie
$\langle O\rangle_{\text{ideal}}$ liczę dokładnie ze statevectora (zamrożona referencja).
Mierzę właśnie tę wielkość, ponieważ: (i) to dokładnie ją redukuje mitygacja (bias
estymatora); (ii) jest ciągła i czuła, w przeciwieństwie do dokładności, która jest progowym
testem znaku i zmienia się skokowo; (iii) jest wprost porównywalna z referencją bezszumową;
(iv) pozwala zlokalizować region załamania ZNE. Dokładność raportuję jako metrykę wtórną.

\paragraph{Fidelity (mirror circuit).} Obwód lustrzany wykonuje $U$, a potem $U^{\dagger}$,
więc stan idealny to dokładnie $|0\rangle^{\otimes n}$, niezależnie od danych wejściowych.
Mierzoną wielkością jest prawdopodobieństwo przeżycia $F = P(|0\rangle^{\otimes n})$.

Wartość $F=1$ oznacza obwód wykonany idealnie. Dolnym
ograniczeniem \emph{nie} jest zero, lecz podłoga stanu maksymalnie zdepolaryzowanego: dla
$n=5$ kubitów losowy rozkład daje $F = 2^{-5} = 0{,}031$. Wynik na poziomie $0{,}03$ oznacza
zatem, że obwód utracił całą informację. Ponieważ obwód lustrzany zawiera $U$ i $U^{\dagger}$,
jego koszt bramkowy jest \emph{dwukrotnością} kosztu obwodu inferencyjnego --- w tabelach
podaję dlatego „mirror CZ”, czyli $2\times$ liczbę CZ obwodu użytkowego. Zaletą tej metryki
jest to, że nie zależy ona od zadania ani od granicy decyzyjnej klasyfikatora: mierzy
wyłącznie jakość wykonania danej struktury routingu na urządzeniu. Praktyczne ograniczenie:
gdy $F$ zbliża się do podłogi $0{,}031$, metryka traci rozdzielczość i nie odróżnia już
obwodów między sobą.

\subsection{Uzasadnienie zbioru danych i konfiguracja \texttt{two\_curves}}

Zbiór \texttt{two\_curves} pochodzi z pakietu \texttt{qml-benchmarks} i towarzyszącego
badania porównawczego modeli QML~\cite{bowles2024}. To syntetyczny zbiór do klasyfikacji
binarnej o kontrolowanej trudności (parametry \texttt{degree}, \texttt{offset}), co pozwala
ustawić nietrywialny, lecz osiągalny pułap klasyczny. Badanie~\cite{bowles2024} pokazuje też,
że przy małych zadaniach klasyczne modele często dorównują kwantowym, a usunięcie splątania
bywa nieszkodliwe --- co uzasadnia traktowanie wzorca splątającego jako obiektu badań, a nie
źródła przewagi.

\paragraph{Zablokowana konfiguracja.} Parametry modelu \texttt{two\_curves} użytego we
wszystkich eksperymentach:

\begin{center}
\begin{tabular}{@{}ll@{}}
\toprule
\textbf{Parametr} & \textbf{Wartość} \\
\midrule
liczba cech $d$              & 8 \\
liczba kubitów $n$           & 5 \\
\texttt{degree}              & 8 \\
kodowanie                    & angle encoding z data re-uploading ($d=8 > n=5$) \\
readout                      & mean value $\frac{1}{n}\sum_q Z_q$ \\
wzorzec splątający (model)   & star \\
liczba warstw $L$ (zablokowana) & 3 \\
dokładność testowa bezszumowa & 0,911 \\
podzbiór ewaluacyjny         & 30 punktów, zbalansowany, seeded \\
\bottomrule
\end{tabular}
\end{center}


\paragraph{Depth sweep.} Model wybrałam wykonując depth sweep po $L \in \{2,3,4,5\}$
i stosując regułę: maksymalna dokładność testowa, remis rozstrzygany mniejszą liczbą bramek
dwukubitowych.

\begin{center}
\begin{tabular}{@{}ccccl@{}}
\toprule
$L$ & dokł. treningowa & dokł. testowa & log. 2q & uwaga \\
\midrule
2 & 0,933 & 0,889 & 8  & \\
3 & 0,933 & \textbf{0,911} & \textbf{12} & \textbf{wybrana} (remis, najmniej 2q) \\
4 & 0,933 & 0,911 & 16 & \\
5 & 0,924 & 0,911 & 20 & \\
\bottomrule
\end{tabular}
\end{center}

\subsection{Uzasadnienie kodowania (angle encoding)}

Stosuję angle encoding: każda cecha wchodzi jako kąt obrotu jednokubitowego, $R_y(x_i)$.
Dla \texttt{two\_curves} angle encoding jest użyty w schemacie \textbf{data re-uploading} ---
ten sam blok kodowania powtarzam w każdej z $L$ warstw, dzięki czemu obwód z $n=5$ kubitami
koduje $d=8 > n$ cech (warunek $d \le n\cdot L$), a funkcja modelu staje się skończonym
szeregiem Fouriera po cechach, co daje większą wyrazistość bez zwiększania liczby
kubitów~\cite{bowles2024}. Dla Banknote używam prostego $R_y$ angle encoding (5 cech na
5 kubitów, bez re-uploadingu). Angle encoding jest tani w bramki i zgodny ze sprzętem, bo
obroty są natywne.

\subsection{Uzasadnienie ansatzu (HEA)}

Używam Hardware-Efficient Ansatz (HEA): każda warstwa to (i) blok angle encoding, (ii) uczone
obroty jednokubitowe, (iii) blok entangling. HEA używa bramek natywnych i lokalnego splątania,
a wzorzec splątający jest regulowanym źródłem bramek dwukubitowych --- ten sam obwód logiczny
przepuszczam przez rosnący koszt routingu. Wybór HEA i traktowanie splątania jako zmiennej
jest spójny z wnioskami benchmarku QML~\cite{bowles2024}.

Poniżej pełne obwody czterech wariantów ansatzu, w kolejności rosnącego kosztu routingu.
Każdy obwód pokazuje jedną warstwę: blok angle encoding $R_y(x)$, warstwę uczonych obrotów
$R(\theta)$ oraz blok entangling złożony z bramek CZ; całość powtarzana jest $L$ razy, po
czym następuje odczyt obserwabli. Warianty różnią się \emph{wyłącznie} blokiem entangling.

\begin{figure}[htbp]
\centering
% --- makra do rysowania obwodu ---
% #1 = etykieta wariantu, #2 = szerokosc bloku entangling, #3 = zawartosc CZ
\newcommand{\ansatzcirc}[3]{%
\begin{tikzpicture}[scale=0.80, baseline,
    gbox/.style={draw,fill=white,minimum width=8mm,minimum height=4.4mm,inner sep=1pt,
                 font=\scriptsize}]
  \def\xend{#2}
  % druty kubitow
  \foreach \q in {0,...,4}{
    \draw (-0.3,-\q*0.72) -- (\xend+1.5,-\q*0.72);
    \node[font=\scriptsize] at (-0.75,-\q*0.72) {$q_\q$};}
  % bloki jednokubitowe
  \foreach \q in {0,...,4}{
    \node[gbox] at (0.55,-\q*0.72) {$R_y(x)$};
    \node[gbox] at (1.85,-\q*0.72) {$R(\theta)$};}
  % separator bloku entangling
  \draw[dashed,gray] (2.5,0.45) -- (2.5,-3.32);
  \draw[dashed,gray] (\xend+0.45,0.45) -- (\xend+0.45,-3.32);
  #3
  % powtorzenie i odczyt
  \node[font=\scriptsize] at (\xend+0.95,0.72) {$\times L$};
  \draw[decorate,decoration={brace,amplitude=3pt}] (\xend+1.5,0.2) -- (\xend+1.5,-3.1);
  \node[font=\scriptsize,anchor=west] at (\xend+1.55,-1.44) {$\langle O\rangle$};
\end{tikzpicture}}
% #1 = x, #2 = qubit a, #3 = qubit b
\newcommand{\czg}[3]{%
  \fill (#1,-#2*0.72) circle (1.7pt);
  \fill (#1,-#3*0.72) circle (1.7pt);
  \draw (#1,-#2*0.72) -- (#1,-#3*0.72);}

\ansatzcirc{star}{5.2}{%
  \czg{2.9}{0}{1}\czg{3.5}{0}{2}\czg{4.1}{0}{3}\czg{4.7}{0}{4}}
\\[4mm]{\small (a) \textbf{star} --- hub $q_0$ splata się z każdym pozostałym kubitem;
4 bramki CZ na warstwę, zgodne z topologią urządzenia (bez SWAP).}

\vspace{5mm}

\ansatzcirc{linear}{5.2}{%
  \czg{2.9}{0}{1}\czg{3.5}{1}{2}\czg{4.1}{2}{3}\czg{4.7}{3}{4}}
\\[4mm]{\small (b) \textbf{linear} --- łańcuch sąsiadów; 4 bramki CZ na warstwę.}

\vspace{5mm}

\ansatzcirc{circular}{5.8}{%
  \czg{2.9}{0}{1}\czg{3.5}{1}{2}\czg{4.1}{2}{3}\czg{4.7}{3}{4}\czg{5.3}{0}{4}}
\\[4mm]{\small (c) \textbf{circular} --- łańcuch domknięty w pierścień; 5 bramek CZ.}

\vspace{5mm}

\ansatzcirc{all-to-all}{8.8}{%
  \czg{2.9}{0}{1}\czg{3.5}{0}{2}\czg{4.1}{0}{3}\czg{4.7}{0}{4}
  \czg{5.3}{1}{2}\czg{5.9}{1}{3}\czg{6.5}{1}{4}
  \czg{7.1}{2}{3}\czg{7.7}{2}{4}\czg{8.3}{3}{4}}
\\[4mm]{\small (d) \textbf{all-to-all} --- każda para kubitów; 10 bramek CZ na warstwę,
najdroższy w routingu.}

\caption{Pełne obwody czterech wariantów ansatzu HEA. Struktura warstwy jest identyczna
(angle encoding $R_y(x)$, uczone obroty $R(\theta)$, blok entangling zaznaczony liniami
przerywanymi), zmienia się wyłącznie blok CZ. Koszt po kompilacji na gwiazdę podaje
tabela~\ref{tab:tc-ladder}.}
\end{figure}

\subsection{Model szumu i kalibracja urządzenia}

Symulator Aer skalibrowałam do 5-kubitowej gwiazdy ODRA (hub QB3 $=$ kubit logiczny 0).
Parametry (kalibracja z 30.07.2026):

\begin{center}
\begin{tabular}{@{}lccccc@{}}
\toprule
& $T_1$ [$\mu$s] & $T_2^{*}$ [$\mu$s] & $T_2^{\text{echo}}$ [$\mu$s] & $p_{1q}$ & $p_{2q}$ \\
\midrule
zakres & 15,6--21,2 & 7,0--10,4 & 11,0--14,7 & $\sim$0,002 & 0,019 \\
\bottomrule
\end{tabular}
\end{center}

\noindent
$p_{ro} \in [0{,}014;\,0{,}026]$; $t_{1q} = 20$~ns, $t_{2q} = 60$~ns.

\wtoku{świeża kalibracja z dnia pomiaru --- do wstawienia; docelowo $\ge 3$ okna
kalibracyjne dla oceny dryfu.}

% =====================================================================
\section{Zbiór danych i ansätze Banknote}
% =====================================================================

Dodatkowo, chcąc sprawdzić \emph{znane} ansätze o różnym dostosowaniu do QPU, wykorzystuję
dane i modele z projektu Banknote (KN-Axion; UCI Banknote, 5 cech, $R_y$ angle encoding,
cechy skalowane do $[-\pi/4, \pi/4]$, etykiety $\{-1,+1\}$). Repozytorium źródłowe:
\url{https://github.com/Kolo-Naukowe-Axion/QC1_Quantum_Banknote_Classifier_on_ODRA_5}.

Stosuję nazewnictwo z pracy źródłowej:

\begin{itemize}
\item \textbf{topology} (w kodzie \texttt{star}) --- natywny, hub QB2, readout $Z$ na
      kubicie 2; kompiluje się na gwiazdę bez SWAP;
\item \textbf{native} (w kodzie \texttt{odra}) --- pierścień $R_Z$/CZ, readout kubit 0;
\item \textbf{baseline} (w kodzie \texttt{simulator}) --- pierścień CRX/CRY, readout
      kubit 0, najdroższy w CZ.
\end{itemize}

Ewaluacja odbywa się na pełnym zbiorze testowym foldu 1 (275 punktów, naturalna
niezbalansowana proporcja 153/122), zgodnie z metodologią
pracy źródłowej. Ze względu na koszt czasu QPU nie wykonuję pełnej walidacji krzyżowej.

% =====================================================================
\section{Wyniki}
% =====================================================================

Sekcja jest zorganizowana zbiorami: najpierw komplet wyników dla \texttt{two\_curves}
(koszt routingu, błąd $\langle O\rangle$ i dokładność, fidelity, wnioski), następnie
analogicznie dla Banknote.

% ---------------------------------------------------------------------
\subsection{\texttt{two\_curves}: routing ladder}
% ---------------------------------------------------------------------

Routing ladder to uporządkowany ciąg wariantów obwodu realizujących \emph{ten sam} obwód
logiczny, ale o rosnącym koszcie CZ po kompilacji --- uzyskany przez zmianę wzorca
entangling. Dzięki temu kontrolowanie zmieniam zmienną niezależną (koszt routingu),
trzymając zadanie stałym. Każdy poziom ma własny klasyfikator wytrenowany przy tej samej
pojemności.

\begin{table}[htbp]
\centering
\caption{Routing ladder \texttt{two\_curves} przy zablokowanym $L=3$.}
\label{tab:tc-ladder}
\begin{tabular}{@{}clccc@{}}
\toprule
Poziom & Wzorzec & log. 2q & CZ po komp. & mirror CZ \\
\midrule
0 & star       & 12 & 12 & 24  \\
1 & linear     & 12 & 27 & 54  \\
2 & circular   & 15 & 36 & 72  \\
3 & all\_to\_all & 30 & 54 & 108 \\
\bottomrule
\end{tabular}
\end{table}

% ---------------------------------------------------------------------
\subsection{\texttt{two\_curves}: błąd $\langle O\rangle$ i dokładność --- symulator vs QPU}
\label{sec:tc-acc}
% ---------------------------------------------------------------------

Parametry pomiaru: 4096 shotów, skale szumu $\{1,3,5\}$, $R=3$ powtórzenia (pasma błędu
i dryf), 30 punktów ewaluacyjnych. Kolumny raportują kolejne szczeble mitygacji:
raw $\rightarrow$ ${+}$REM $\rightarrow$ ${+}$ZNE (Richardson).

\begin{table}[htbp]
\centering
\caption{\texttt{two\_curves}: błąd $\langle O\rangle$ i dokładność. Symulator (DD${+}$REM)
oraz rzeczywiste ODRA (wiersze DD-off), $R=3$.}
\begin{tabular}{@{}lc ccc ccc@{}}
\toprule
& & \multicolumn{3}{c}{\textbf{błąd $\langle O\rangle$}} & \multicolumn{3}{c}{\textbf{accuracy}} \\
\cmidrule(lr){3-5}\cmidrule(lr){6-8}
Wzorzec & CZ & raw & ${+}$REM & ${+}$ZNE & raw & ${+}$REM & ${+}$ZNE \\
\midrule
\multicolumn{8}{@{}l}{\emph{Symulator}} \\
star        & 12 & 0,192 & 0,162 & 0,024 & 0,867 & 0,867 & 0,867 \\
linear      & 27 & 0,406 & 0,378 & 0,240 & 0,878 & 0,900 & 0,900 \\
circular    & 36 & 0,402 & 0,380 & 0,304 & 0,900 & 0,900 & 0,900 \\
all\_to\_all & 54 & 0,456 & 0,436 & 0,351 & 0,889 & 0,900 & 0,900 \\
\midrule
\multicolumn{8}{@{}l}{\emph{Rzeczywiste ODRA (QPU), DD-off}} \\
star        & 12 & 0,112 & 0,088 & \textbf{0,024} & 0,900 & 0,900 & 0,867 \\
linear      & 27 & 0,281 & 0,267 & 0,115 & 0,922 & 0,922 & 0,900 \\
circular    & 36 & 0,299 & 0,292 & 0,176 & 0,900 & 0,900 & 0,900 \\
all\_to\_all & 54 & 0,357 & 0,349 & 0,239 & 0,878 & 0,878 & 0,867 \\
\midrule
\multicolumn{8}{@{}l}{\emph{Rzeczywiste ODRA (QPU), DD-on --- ablacja}} \\
star        & 12 & 0,159 & 0,138 & 0,029 & 0,900 & 0,900 & 0,867 \\
linear      & 27 & 0,375 & 0,365 & 0,209 & 0,900 & 0,900 & 0,900 \\
circular    & 36 & 0,382 & 0,378 & 0,292 & 0,911 & 0,900 & 0,911 \\
all\_to\_all & 54 & 0,458 & 0,454 & 0,373 & 0,900 & 0,900 & 0,911 \\
\bottomrule
\end{tabular}
\end{table}

\begin{figure}[htbp]
\centering
\includegraphics[width=0.9\textwidth]{figures/fig_two_curves_qpu_vs_sim.png}
\caption{\texttt{two\_curves}: $|\langle O\rangle$ err$|$ vs CZ, symulator vs rzeczywiste
ODRA, szczeble raw / ${+}$REM / ${+}$ZNE.}
\end{figure}

\paragraph{Wnioski.}
\begin{enumerate}
\item \textbf{Błąd rośnie monotonicznie z kosztem CZ} --- i na symulatorze, i na sprzęcie.
      Routing degraduje wynik, co potwierdza wybór CZ jako właściwej zmiennej niezależnej.
\item \textbf{ZNE działa i jest silne.} Na QPU przy star błąd spada z $0{,}112$ do
      $0{,}024$, czyli \textbf{4,6-krotnie}. Odzysk maleje z kosztem: przy all\_to\_all
      z $0{,}357$ tylko do $0{,}239$.
\item \textbf{Dokładność nie reaguje.} Dla każdego z czterech wzorców entangling dokładność
      pozostaje w przedziale $0{,}867$--$0{,}922$, niezależnie od szczebla mitygacji. Jest to
      bezpośrednia ilustracja ram z sekcji 2.1: predykcja to $\mathrm{sgn}\langle O\rangle$,
      a więc próg, który sam z siebie tłumi szum. Punkty leżą daleko od granicy decyzyjnej,
      więc redukcja biasu nie zmienia znaku --- to błąd $\langle O\rangle$ niesie tu sygnał.
\item \textbf{DD pogarsza wynik na sprzęcie.} Przy każdym wzorcu wiersz DD-on ma wyższy błąd
      niż DD-off (star: $0{,}159$ vs $0{,}112$; all\_to\_all: $0{,}458$ vs $0{,}357$).
\end{enumerate}

% ---------------------------------------------------------------------
\subsection{\texttt{two\_curves}: mirror fidelity}
% ---------------------------------------------------------------------

Pomiar wykonany na rzeczywistym QPU (8192 shotów, 5 reprezentatywnych wejść) oraz
w symulatorze przy identycznym protokole.

\begin{table}[htbp]
\centering
\caption{\texttt{two\_curves}: mirror fidelity $P(|0\rangle^{5})$. Podłoga depolaryzacyjna
wynosi $2^{-5} = 0{,}031$.}
\begin{tabular}{@{}lc cc cc c@{}}
\toprule
& & \multicolumn{2}{c}{\textbf{QPU}} & \multicolumn{2}{c}{\textbf{symulator}} & \\
\cmidrule(lr){3-4}\cmidrule(lr){5-6}
Wzorzec & mirror CZ & DD-off & DD-on & DD-off & DD-on & QPU/sim \\
\midrule
star        & 24  & \textbf{0,549} & 0,389 & 0,330 & 0,367 & 1,67$\times$ \\
linear      & 54  & 0,211 & 0,120 & 0,085 & 0,102 & 2,48$\times$ \\
circular    & 72  & 0,118 & 0,060 & 0,051 & 0,062 & 2,31$\times$ \\
all\_to\_all & 108 & 0,094 & 0,048 & 0,054 & 0,059 & 1,73$\times$ \\
\bottomrule
\end{tabular}
\end{table}

\paragraph{Dlaczego te wartości są tak niskie.} Bezwzględny poziom fidelity może wydawać
się zaskakująco mały, jest jednak dokładnie taki, jakiego należy oczekiwać, i nie wynika
z błędu pomiaru. Składają się na to trzy niezależne czynniki.

Po pierwsze, $P(|0\rangle^{5})$ jest kryterium bardzo surowym: zdarzenie jest zaliczone
tylko wtedy, gdy \emph{wszystkie pięć} kubitów zostanie odczytanych poprawnie
jednocześnie. Pojedynczy błąd na dowolnym kubicie niszczy cały wynik, więc wielkość ta
spada znacznie szybciej niż jednokubitowa wartość oczekiwana $\langle Z\rangle$, do której
przywykliśmy w pozostałych sekcjach.

Po drugie, obwód lustrzany z definicji \emph{podwaja} koszt bramkowy: wariant star używa
w inferencji 12 bramek CZ, ale jego wersja lustrzana już 24. Mierzymy zatem obwód dwa razy
dłuższy niż ten, który realizuje klasyfikację.

Po trzecie, degradacja jest wykładnicza w liczbie bramek. Dopasowany model urządzenia
(sekcja 5.7) daje efektywnie $0{,}86\%$ błędu na bramkę CZ i $0{,}613\%$ na bramkę
jednokubitową. Dla wariantu star (24 CZ i 80 bramek jednokubitowych) przewiduje to
$F \approx \exp(-0{,}0086\cdot 24 - 0{,}00613\cdot 80) \approx 0{,}50$, wobec zmierzonych
$0{,}549$. Zgodność jest bardzo dobra, co potwierdza, że pomiar jest spójny z niezależnie
wyznaczonymi parametrami błędu.

\paragraph{Wnioski.}
\begin{enumerate}
\item \textbf{Fidelity spada monotonicznie z kosztem routingu} --- od $0{,}549$ przy star do
      $0{,}094$ przy all\_to\_all. Trend jest niezależny od zadania, więc potwierdza, że
      degradacja jest własnością struktury obwodu, a nie klasyfikatora.
\item \textbf{QPU jest o 1,7--2,5$\times$ lepsze niż nasz symulator} przy każdym poziomie.
      Kalibrowany model szumu jest systematycznie pesymistyczny, co jest spójne z wnioskiem
      z pomiaru błędu $\langle O\rangle$.
\item \textbf{DD pogarsza fidelity na sprzęcie, ale poprawia ją w symulatorze.} Na QPU każda konfiguracja z DD wypada gorzej
      (star: $0{,}389$ wobec $0{,}549$), podczas gdy w symulatorze DD nieznacznie pomaga
      ($0{,}367$ wobec $0{,}330$). Przyczyna jest jasna: w symulatorze modeluję DD jako
      poprawę $T_2^{*}\to T_2^{\text{echo}}$, więc z konstrukcji może ono wyłącznie pomóc.
      Rzeczywiste impulsy DD wnoszą natomiast własny błąd koherentny, którego model nie
      zawiera. Wniosek praktyczny: \emph{o skuteczności DD nie wolno wnioskować z symulacji},
      trzeba ją zmierzyć na urządzeniu --- i na IQM SPARKu wypada ona negatywnie.
\item \textbf{Najdroższe warianty wyczerpują metrykę.} Wartości $0{,}118$ i $0{,}094$ są już
      tylko 3--4$\times$ powyżej podłogi $0{,}031$, więc mirror fidelity traci tam
      rozdzielczość i nie należy nadinterpretować różnicy między circular a all\_to\_all.
\end{enumerate}

% ---------------------------------------------------------------------
\subsection{Banknote: baseline bezszumowy i routing ladder kosztu}
% ---------------------------------------------------------------------

\begin{table}[htbp]
\centering
\caption{Banknote: dokładność bezszumowa (średnia z 5 foldów) i koszt CZ na gwieździe.}
\begin{tabular}{@{}llccc@{}}
\toprule
Ansatz (nazwa w kodzie) & głęb. & CZ & głęb. po komp. & dokł. bezszum. \\
\midrule
topology (\texttt{star})      & 2/4/6 & 8/16/24  & 14/26/38   & 0,856 / 0,889 / 0,885 \\
native (\texttt{odra})        & 2/4/6 & 13/35/57 & 24/66/108  & 0,871 / 0,876 / 0,886 \\
baseline (\texttt{simulator}) & 2/4/6 & 20/52/87 & 56/140/231 & 0,891 / 0,880 / 0,883 \\
\bottomrule
\end{tabular}
\end{table}

\paragraph{Wniosek.} Trzy ansätze mają niemal identyczną dokładność bezszumową (około
$0{,}88$), ale różnią się kosztem CZ nawet 3,6-krotnie (24 vs 87 przy głębokości 6). Jest to
więc czysta routing ladder: przy ustalonym zadaniu i ustalonej jakości modelu zmienia się
niemal wyłącznie koszt kompilacji.

% ---------------------------------------------------------------------
\subsection{Banknote: badanie ablacyjne w symulatorze}
\label{sec:bank-sim}
% ---------------------------------------------------------------------

\begin{table}[htbp]
\centering
\caption{Banknote (symulator): błąd $\langle O\rangle$ raw $\to$ ZNE oraz dokładność
raw $\to$ ZNE.}
\begin{tabular}{@{}llc ccc c@{}}
\toprule
Ansatz & głęb. & CZ & err raw & err ${+}$DD${+}$REM & err ZNE & dokł. (raw $\to$ ZNE) \\
\midrule
topology & 2 & 8  & 0,104 & 0,090 & 0,030 & 0,78 $\to$ 0,78 \\
native   & 2 & 13 & 0,088 & 0,076 & 0,026 & 0,78 $\to$ 0,76 \\
baseline & 2 & 20 & 0,144 & 0,135 & 0,056 & 0,82 $\to$ 0,78 \\
\midrule
topology & 4 & 16 & 0,155 & 0,144 & 0,056 & 0,82 $\to$ 0,80 \\
native   & 4 & 35 & 0,215 & 0,210 & 0,130 & 0,82 $\to$ 0,78 \\
baseline & 4 & 55 & 0,256 & 0,248 & 0,175 & 0,80 $\to$ 0,80 \\
\midrule
topology & 6 & 24 & 0,187 & 0,174 & 0,091 & 0,82 $\to$ 0,76 \\
native   & 6 & 60 & 0,251 & 0,249 & 0,212 & 0,80 $\to$ 0,78 \\
baseline & 6 & 87 & 0,294 & 0,282 & 0,238 & \textbf{0,60 $\to$ 0,80} \\
\bottomrule
\end{tabular}
\end{table}

\begin{figure}[htbp]
\centering
\includegraphics[width=0.9\textwidth]{figures/fig_banknote_ladder.png}
\caption{Banknote (symulator): błąd $\langle O\rangle$ (lewy panel) i dokładność (prawy)
vs CZ, zbiorczo dla głębokości $\{2,4,6\}$.}
\end{figure}

\paragraph{Wnioski.} Najważniejszy wynik tej sekcji: \textbf{ZNE potrafi odzyskać
dokładność klasyfikatora}. Dla wariantu baseline przy głębokości 6 (87 CZ) mitygacja
podnosi dokładność z $0{,}60$ do $0{,}80$, czyli praktycznie do poziomu bezszumowego --- szum
zdążył tam już przerzucić znak $\langle O\rangle$ dla części punktów, a ZNE cofa ten efekt.

Poza tym jednym przypadkiem krzywa ZNE leży wprawdzie poniżej raw i DD${+}$REM dla każdej
konfiguracji, ale różnica maleje wraz z kosztem CZ, więc odzysk słabnie. Potwierdza to
obserwację z \texttt{two\_curves}: dokładność reaguje na mitygację dopiero wtedy, gdy szum
jest na tyle duży, że zaczyna zmieniać znak wartości oczekiwanej.

% ---------------------------------------------------------------------
\subsection{Banknote: rzeczywiste QPU}
\label{sec:bank-qpu}
% ---------------------------------------------------------------------

Wynik główny tej sekcji pochodzi z \textbf{pełnego foldu testowego (275 punktów, fold~1,
naturalna proporcja klas 153/122, podłoga większości $0{,}56$)} --- czyli dokładnie ze
zbioru ewaluacyjnego pracy źródłowej --- wykonanego na rzeczywistym QPU dla wszystkich
trzech ansätze i trzech głębokości. Metryki: błąd $\langle O\rangle$ (raw${}\to{}$ZNE),
dokładność (ideał / raw / ZNE) oraz macro-F1 (raw / ZNE). DD wyłączone (uzasadnienie:
wynik negatywny DD, sekcja poniżej i 5.7).

\begin{table}[htbp]
\centering
\caption{Banknote na ODRA --- pełny fold (275 pkt), DD-off. Dla każdego ansatzu podano błąd
$\langle O\rangle$ raw${}\to{}$ZNE, dokładność (ideał bezszumowy / raw / ZNE) oraz macro-F1
(raw / ZNE). ZNE = REM${+}$ekstrapolacja Richardsona.}
\label{tab:bank-qpu-full}
\begin{tabular}{@{}llc cc ccc cc@{}}
\toprule
& & & \multicolumn{2}{c}{err $\langle O\rangle$} & \multicolumn{3}{c}{dokładność} & \multicolumn{2}{c}{macro-F1} \\
\cmidrule(lr){4-5}\cmidrule(lr){6-8}\cmidrule(lr){9-10}
Ansatz & głęb. & CZ & raw & ZNE & ideał & raw & ZNE & raw & ZNE \\
\midrule
topology & 2 & 8  & 0,100 & 0,039 & 0,822 & 0,876 & 0,815 & 0,873 & 0,801 \\
native   & 2 & 13 & 0,067 & 0,033 & 0,836 & 0,873 & 0,840 & 0,871 & 0,834 \\
baseline & 2 & 20 & 0,152 & 0,063 & 0,858 & 0,887 & 0,876 & 0,889 & 0,874 \\
\midrule
topology & 4 & 16 & 0,111 & 0,030 & 0,840 & 0,851 & 0,829 & 0,843 & 0,815 \\
native   & 4 & 35 & 0,161 & 0,050 & 0,833 & 0,858 & 0,829 & 0,857 & 0,815 \\
baseline & 4 & 52 & 0,280 & 0,226 & 0,833 & 0,706 & \textbf{0,767} & 0,699 & \textbf{0,767} \\
\midrule
topology & 6 & 24 & 0,142 & 0,057 & 0,836 & 0,800 & 0,804 & 0,778 & 0,783 \\
native   & 6 & 57 & 0,193 & 0,127 & 0,833 & 0,800 & 0,826 & 0,781 & 0,815 \\
baseline & 6 & 87 & 0,291 & 0,279 & 0,833 & 0,673 & \textbf{0,727} & 0,663 & \textbf{0,724} \\
\bottomrule
\end{tabular}
\end{table}

\begin{figure}[htbp]
\centering
\includegraphics[width=0.95\textwidth]{figures/fig_banknote_qpu.png}
\caption{Banknote na rzeczywistym ODRA (pełny fold, 275 pkt, DD-off): błąd $\langle O\rangle$
(lewy panel) oraz dokładność (prawy panel) wobec transpilowanego kosztu CZ, dla drabiny
routingu topology${}<{}$native${}<{}$baseline i głębokości $\{2,4,6\}$, wzdłuż potoku
raw${}\to{}$+REM${}\to{}$+ZNE. Analogiczny do rysunku symulatorowego z sekcji~\ref{sec:bank-sim}.
Dane z tabeli~\ref{tab:bank-qpu-full}.}
\label{fig:bank-qpu}
\end{figure}

\paragraph{Główny wynik sprzętowy --- odzysk dokładności i słabnięcie ZNE z kosztem CZ.}
Poniższe wnioski opisują pomiar pojedynczego foldu (fold~1, tab.~\ref{tab:bank-qpu-full});
wersję odporną, uśrednioną po pięciu foldach, podaje tab.~\ref{tab:bank-qpu-cv} niżej (baseline
jest silnie zależny od foldu, więc liczby dla 87~CZ należy czytać łącznie z nią). Dane na
pełnym foldzie układają się w spójny obraz sterowany kosztem routingu.

\begin{enumerate}
\item \textbf{Współczynnik redukcji błędu ZNE maleje monotonicznie z CZ.} Przy niskim koszcie
      ZNE tnie błąd $\langle O\rangle$ silnie (topology, 16 CZ: $0{,}111\to0{,}030$, czyli
      $3{,}7\times$), przy wysokim coraz słabiej, aż do granicy skuteczności: baseline przy
      87 CZ daje już tylko $0{,}291\to0{,}279$ (współczynnik $1{,}04$, mitygacja praktycznie
      nieskuteczna). Jest to najczystsza sprzętowa ilustracja słabnięcia ZNE z kosztem CZ.
\item \textbf{Dokładność reaguje wyłącznie tam, gdzie szum zdążył przerzucić znak
      $\langle O\rangle$.} Dla tanich ansätze (topology, native, oraz baseline przy
      głębokości 2) dokładność surowa jest już bliska ideału, więc ZNE nie ma czego
      odzyskiwać i bywa nieznacznie niższa (nasycenie, dyskusja w sekcji 2.1). Odzysk pojawia
      się dopiero przy baseline i dużym koszcie: przy 52 CZ dokładność rośnie
      $0{,}706\to0{,}767$ (macro-F1 $0{,}699\to0{,}767$), a przy 87 CZ $0{,}673\to0{,}727$
      (macro-F1 $0{,}663\to0{,}724$) --- odzysk rzędu $34$--$48\%$ luki do ideału, mierzony na
      rzeczywistym sprzęcie.
\item \textbf{Macro-F1 podąża za dokładnością, więc odzysk jest zrównoważony klasowo.}
      W żadnym wierszu DD-off macro-F1 nie odrywa się od dokładności (np. baseline 87 CZ:
      dokł. $0{,}673$ / F1 $0{,}663$), czyli klasyfikator nie zapada się w stronę klasy
      większościowej --- mimo naturalnej niezbalansowanej proporcji foldu.
\end{enumerate}


\paragraph{DD na sprzęcie --- ablacja (podzbiór 40 pkt, głębokość 6).}
Poniższa tabela (te same obwody, DD włączone) dokumentuje wynik negatywny DD, który na
pełnym foldzie nie był mierzony.

\begin{table}[htbp]
\centering
\caption{Banknote na ODRA, głębokość 6 --- ablacja DD (podzbiór zbalansowany 40 pkt,
te same obwody, DD włączone).}
\begin{tabular}{@{}lc cc cc@{}}
\toprule
Ansatz & CZ & err raw & err ZNE & dokł. raw & dokł. ZNE \\
\midrule
topology & 24 & 0,183 & 0,140 & 0,517 & 0,650 \\
native   & 60 & 0,236 & 0,198 & 0,483 & 0,642 \\
baseline & 84 & 0,232 & 0,198 & 0,525 & 0,558 \\
\bottomrule
\end{tabular}
\end{table}

\noindent Przy każdym ansatzu włączenie DD wyraźnie obniża dokładność względem DD-off
(topology: $0{,}733\to0{,}517$; native: $0{,}775\to0{,}483$), co potwierdza wynik negatywny
DD z sekcji 5.7 i uzasadnia raportowanie potoku jako REM${+}$ZNE.

\paragraph{Porównanie z pracą źródłową (raw wobec raw) i walidacja krzyżowa 5-fold.}
Praca źródłowa mierzy dokładność i macro-F1 \emph{na rzeczywistym QPU} (IQM Spark), na
zamrożonych wagach, \emph{bez mitygacji} (raw), uśredniając po pięciu foldach walidacji
krzyżowej. Właściwym punktem odniesienia jest zatem nasza kolumna \emph{raw} (bez REM/ZNE);
REM i ZNE stanowią \emph{dodatkowy} poziom odzysku ponad tym, co raportuje praca. Aby
porównać w identycznych warunkach, powtórzyłam pomiar na wszystkich pięciu foldach przy
głębokości 6 (native i baseline, DD-off, 275 pkt/fold, 4096 shots); wynik zapisany w
\texttt{banknote\_hardware\_cv\_d6.csv}.

\begin{table}[htbp]
\centering
\caption{Banknote L=6 na ODRA --- walidacja krzyżowa 5-fold (275 pkt/fold, DD-off, 4096 shots),
średnia${}\pm{}$odchylenie międzyfoldowe. Kolumna \emph{raw} jest bezpośrednio porównywalna z
pracą źródłową (QPU, bez mitygacji); ZNE ${}={}$ REM${+}$ekstrapolacja Richardsona.}
\label{tab:bank-qpu-cv}
\begin{tabular}{@{}lc cc cc c@{}}
\toprule
& & \multicolumn{2}{c}{dokładność} & \multicolumn{2}{c}{macro-F1} & praca (QPU) \\
\cmidrule(lr){3-4}\cmidrule(lr){5-6}
Ansatz & CZ & raw & ZNE & raw & ZNE & dokł. / F1 \\
\midrule
native (odra)  & 57--60 & $0{,}809\pm0{,}042$ & $0{,}867\pm0{,}030$ & $0{,}796\pm0{,}049$ & $0{,}862\pm0{,}033$ & $0{,}590$ / $0{,}142$ \\
baseline (sim) & 84--87 & $0{,}626\pm0{,}091$ & $0{,}643\pm0{,}097$ & $0{,}601\pm0{,}107$ & $0{,}624\pm0{,}108$ & $0{,}565$ / $0{,}044$ \\
\bottomrule
\end{tabular}
\end{table}

\noindent Wnioski z walidacji krzyżowej:
\begin{enumerate}
\item \textbf{Nie odtwarzam zapaści pierścieni z pracy.} Praca raportuje przy L=6 zapaść obu
      pierścieni do podłogi klasy większościowej (native dokł. $0{,}590$ / F1 $0{,}142$;
      baseline $0{,}565$ / F1 $0{,}044$). U nas native pozostaje \emph{zdrowy}
      (raw $0{,}809$ / F1 $0{,}796$), a baseline degraduje się w stronę podłogi dokładności,
      ale zachowuje F1${}\approx0{,}60$ --- czyli \emph{nie} zapada się w stronę klasy
      większościowej. Nawet nasz surowy, niemitygowany wynik jest wyraźnie powyżej pracy.
\item \textbf{Różnica nie pochodzi z naszej mitygacji.} Wagi, ansätze, odczyt (q0), kodowanie,
      transpilacja (star coupling, \texttt{optimization\_level=1}) i rodzina urządzenia są
      identyczne z pracą, a nasz \emph{raw} już ją przewyższa. Jedyną pozostałą zmienną jest
      stan urządzenia: kalibracja, dryf, crosstalk oraz to, na które fizyczne kubity trafił
      layout. Bezpośredni dowód dryfu mam we własnych danych: ten sam fold~1 baseline L=6 dał
      raw $0{,}673$ w pomiarze z 1~sierpnia (tab.~\ref{tab:bank-qpu-full}) i $0{,}607$ w
      pomiarze 5-foldowym z 3~sierpnia --- różnica $0{,}066$ na identycznych obwodach.
\item \textbf{To wzmacnia tezę pracy.} Praca argumentuje, że zachowanie on-device jest zależne
      od kalibracji i nie transferuje z symulacji; nasz wynik dokłada mocniejszą wersję: sama
      \emph{zapaść} pierścieni nie jest powtarzalna między kalibracjami tego samego urządzenia.
\item \textbf{ZNE odzyskuje ponad raw także tutaj, ale słabnie z kosztem CZ.} Native (57 CZ):
      raw $0{,}809\to$ ZNE $0{,}867$ (błąd $\langle O\rangle$ $0{,}190\to0{,}119$, $1{,}6\times$).
      Baseline (87 CZ): raw $0{,}626\to$ ZNE $0{,}643$ (błąd $0{,}278\to0{,}284$) --- przy tym
      koszcie redukcja błędu \emph{zanika} i uśredniona po foldach potrafi się odwrócić, co
      potwierdza górną granicę użyteczności ZNE z sekcji~\ref{sec:takeaways}.
\item \textbf{Duża rozbieżność międzyfoldowa baseline} (odch.\ std dokł.\ $\pm0{,}091$ wobec
      $\pm0{,}042$ dla native) pokazuje, że pojedynczy fold jest tu zawodny --- to właśnie
      uzasadniało konieczność pełnej walidacji 5-fold.
\end{enumerate}

% ---------------------------------------------------------------------
\subsection{Banknote: mirror fidelity na QPU}
% ---------------------------------------------------------------------

\begin{table}[htbp]
\centering
\caption{Banknote: mirror fidelity $P(|0\rangle^{5})$ na rzeczywistym QPU (8192 shotów).}
\begin{tabular}{@{}lccc cc@{}}
\toprule
Ansatz & głęb. & mirror CZ & DD-off & DD-on & DD-on / DD-off \\
\midrule
topology & 2 & 16  & 0,619 & 0,425 & 0,69 \\
topology & 4 & 32  & 0,501 & 0,280 & 0,56 \\
topology & 6 & 48  & 0,245 & 0,137 & 0,56 \\
\midrule
native   & 2 & 26  & 0,596 & 0,386 & 0,65 \\
native   & 4 & 70  & 0,226 & 0,109 & 0,48 \\
native   & 6 & 120 & 0,102 & 0,043 & 0,42 \\
\midrule
baseline & 2 & 40  & 0,345 & 0,225 & 0,65 \\
baseline & 4 & 110 & 0,128 & 0,057 & 0,45 \\
baseline & 6 & 174 & 0,051 & 0,032 & 0,63 \\
\bottomrule
\end{tabular}
\end{table}

\begin{figure}[htbp]
\centering
\includegraphics[width=0.85\textwidth]{figures/fig_mirror_fidelity.png}
\caption{Mirror-circuit fidelity vs koszt routingu.}
\end{figure}

\paragraph{Wnioski.}
\begin{enumerate}
\item Fidelity spada monotonicznie z mirror CZ w obrębie każdego ansatzu i porządkuje ansätze
      dokładnie tak, jak przewiduje koszt routingu: topology $>$ native $>$ baseline.
\item \textbf{DD pogarsza fidelity we wszystkich dziewięciu konfiguracjach}, zachowując
      jedynie 42--69\% wartości DD-off, przy czym kara rośnie z głębokością obwodu.
\item Baseline przy głębokości 6 osiąga $0{,}051$, czyli praktycznie podłogę
      depolaryzacyjną $0{,}031$ --- obwód traci niemal całą informację. To niezależne
      potwierdzenie, że słabnięcie ZNE przy wysokim koszcie CZ zbiega się z reżimem, w którym
      obwód przestaje w ogóle przenosić sygnał.
\item Łącząc wszystkie 13 punktów pomiarowych QPU z obu badań (\texttt{two\_curves}
      i Banknote), jednolity dwuparametrowy model
      $F = (1-2^{-n})\exp(-a\,n_{CZ} - b\,n_{1q}) + 2^{-n}$ odtwarza pomiary z medianą błędu
      względnego $10{,}5\%$ przy efektywnym $0{,}86\%$ na bramkę CZ i $0{,}613\%$ na bramkę
      jednokubitową. Ponieważ obwody zawierają 2--4$\times$ więcej bramek jednokubitowych niż
      CZ, \emph{to bramki jednokubitowe odpowiadają za około 60\% budżetu błędu}. Sam koszt
      CZ jest więc dobrym, ale niepełnym predyktorem jakości na sprzęcie.
\end{enumerate}

% =====================================================================
\section{Najważniejsze obserwacje z Miesiąca 1}
\label{sec:takeaways}
% =====================================================================

Poniżej zbieram najważniejsze, zweryfikowane wnioski z całego miesiąca. Każdy jest poparty
konkretnymi liczbami z sekcji wyników; podaję je w formie skrótowej jako podsumowanie
wykonawcze.

\begin{enumerate}
\item \textbf{ZNE działa, a jego siła słabnie przewidywalnie z kosztem CZ.} Na rzeczywistym
      QPU (\texttt{two\_curves}, DD-off) współczynnik redukcji błędu $\langle O\rangle$ spada
      wraz z routingiem: star (12 CZ) $0{,}112\to0{,}024$ ($4{,}6\times$), linear (27 CZ)
      $0{,}281\to0{,}115$ ($2{,}4\times$), circular (36 CZ) $0{,}299\to0{,}176$ ($1{,}7\times$),
      all\_to\_all (54 CZ) $0{,}357\to0{,}239$ ($1{,}5\times$). Trend jest monotoniczny i
      powtarza się na symulatorze.

\item \textbf{Dokładność jest metryką „głuchą”, dopóki szum nie przerzuci znaku
      $\langle O\rangle$.} W \texttt{two\_curves} dokładność trzyma się w przedziale
      $0{,}87$--$0{,}92$ mimo $15\times$ zmiany błędu $\langle O\rangle$. Mitygacja zaczyna
      poprawiać dokładność dopiero \emph{powyżej progu kosztu}, gdy szum przerzucił znak
      $\langle O\rangle$ dla części punktów --- na QPU (Banknote, baseline) próg leży około
      50 CZ: przy 52 CZ dokładność rośnie $0{,}706\to0{,}767$, a przy 87 CZ
      $0{,}673\to0{,}727$. Poniżej progu dokładność surowa jest już przy suficie i ZNE jej nie
      podnosi. To bezpośrednia ilustracja ram bias/wariancja z sekcji 2.1.

\item \textbf{Dynamical Decoupling to wynik negatywny na tym urządzeniu --- i symulator tego
      nie widzi.} Natywne \texttt{STANDARD\_DD} pogorszyło każdą mierzoną wielkość w czterech
      niezależnych eksperymentach:
      \begin{center}
      \begin{tabular}{@{}lcc@{}}
      \toprule
      eksperyment & DD-off & DD-on \\
      \midrule
      \texttt{two\_curves} star, błąd $\langle O\rangle$        & 0,112 & 0,159 \\
      \texttt{two\_curves} star, mirror fidelity                 & 0,549 & 0,389 \\
      Banknote d6 star, dokładność QPU                           & 0,733 & 0,517 \\
      Banknote d6 star, mirror fidelity                          & 0,245 & 0,137 \\
      \bottomrule
      \end{tabular}
      \end{center}
      Kalibrowany symulator przewiduje, że DD \emph{pomaga}, bo modeluje je wyłącznie jako
      poprawę $T_2^{*}\to T_2^{\text{echo}}$. Wniosek metodologiczny: skuteczność DD trzeba
      mierzyć na urządzeniu.

\item \textbf{Kalibrowany symulator jest konserwatywnym, pesymistycznym przybliżeniem.}
      Każdy punkt mirror fidelity na sprzęcie przewyższa predykcję symulatora o czynnik od
      $1{,}7$ do $2{,}5$ (\texttt{two\_curves}: star $0{,}549$ wobec $0{,}330$, linear
      $0{,}211$ wobec $0{,}085$). Trening względem symulatora daje więc bezpieczne dolne
      ograniczenie, a nie zawyżoną prognozę.

\item \textbf{Mirror fidelity opisuje jeden dwuparametrowy model per-bramka.} Dopasowanie do
      13 punktów QPU z obu badań daje $\approx0{,}86\%$ błędu na bramkę CZ i $\approx0{,}61\%$
      na bramkę jednokubitową (mediana błędu względnego $10{,}5\%$); dla star model przewiduje
      $F\approx0{,}50$ wobec zmierzonych $0{,}549$. Ponieważ mirror fidelity jest niezależna od
      zadania ($U U^{\dagger}$), degradacja jest własnością struktury obwodu, nie
      klasyfikatora.

\item \textbf{Koszt CZ to właściwa, lecz niepełna zmienna niezależna.} Błąd rośnie
      monotonicznie z CZ w dwóch różnych zbiorach i niezależnie od ansatzu, co uzasadnia całą
      konstrukcję routing ladder. Zarazem obwody mają $2$--$4\times$ więcej bramek
      jednokubitowych niż CZ, więc te pierwsze odpowiadają za około $60\%$ budżetu błędu ---
      sam CZ jest dobrym, ale niepełnym predyktorem.

\item \textbf{ZNE ma dwa różne profile zależności od kosztu CZ --- i trzeba je rozróżnić.}
      \begin{itemize}
      \item \emph{Redukcja błędu $\langle O\rangle$ słabnie monotonicznie z CZ.} Jest
            najsilniejsza przy niskim koszcie i zanika przy wysokim: współczynnik spada z
            $4{,}6\times$ (12 CZ) do $1{,}04\times$ (87 CZ). Dla tej metryki nie ma „okna" ---
            ZNE jest po prostu najlepsze przy tanich obwodach.
      \item \emph{Odzysk dokładności ma natomiast charakter okna} („wystarczająco dużo szumu,
            ale nie za dużo"). Pojawia się dopiero powyżej pewnego progu kosztu (tu $\gtrsim50$
            CZ), gdy szum zdąży obniżyć dokładność surową poniżej ideału i jest co odzyskiwać
            (baseline 52 CZ: $0{,}706\to0{,}767$; 87 CZ: $0{,}673\to0{,}727$). Poniżej progu
            dokładność surowa jest już przy suficie, więc ZNE nic nie wnosi i bywa nieznacznie
            niższa (dodana wariancja).
      \end{itemize}
      Dolna krawędź okna jest zmierzona. Górną krawędź lokalizuje walidacja 5-foldowa: przy
      87 CZ (baseline L=6) uśredniona po foldach redukcja błędu $\langle O\rangle$
      \emph{zanika i potrafi się odwrócić} ($0{,}278\to0{,}284$), a przyrost dokładności
      ($0{,}626\to0{,}643$) mieści się w rozrzucie międzyfoldowym --- czyli w okolicach
      $85$--$90$ CZ ZNE przestaje być opłacalne. Dokładne wyznaczenie progu na innych
      urządzeniach jest zadaniem Miesiąca~2. Dodatkowo wariant wykładniczy ekstrapolacji staje
      się niestabilny przy wysokim koszcie.

\item \textbf{Mirror fidelity ma granicę rozdzielczości przy podłodze depolaryzacyjnej}
      $2^{-5}=0{,}031$. Dla all\_to\_all (\texttt{two\_curves}) wynosi $0{,}094$ --- tylko
      $\sim3\times$ powyżej podłogi --- więc różnic circular vs all\_to\_all nie należy
      nadinterpretowywać.

\item \textbf{Nie odtwarzamy zapaści pierścieni z pracy źródłowej --- i jest to efekt stanu
      urządzenia, nie mitygacji.} Praca raportuje na QPU przy L=6 zapaść obu pierścieni do
      podłogi klasy większościowej (native dokł.\ $0{,}590$ / F1 $0{,}142$; baseline $0{,}565$
      / F1 $0{,}044$). W naszej walidacji 5-foldowej --- przy identycznych wagach, ansätzach,
      odczycie i transpilacji --- native pozostaje zdrowy (\emph{raw} $0{,}809$ / F1 $0{,}796$),
      a baseline degraduje się w stronę podłogi dokładności, lecz \emph{nie} zapada klasowo
      (F1${}\approx0{,}60$). Ponieważ nawet nasz surowy wynik przewyższa pracę, różnicę
      przypisuję dryfowi kalibracji (bezpośredni dowód: ten sam fold~1 baseline dał raw
      $0{,}673$ i $0{,}607$ w dwóch pomiarach z różnych dni). Macro-F1 podąża przy tym za
      dokładnością (brak zapaści klasowej w wierszach DD-off), a przy DD-on spada do
      $\sim0{,}52$ (blisko podłogi). Szczegóły w sekcji~\ref{sec:bank-qpu} (tab.~\ref{tab:bank-qpu-cv}).
\end{enumerate}

\paragraph{Podsumowanie} QEM (REM${+}$ZNE) mierzalnie odzyskuje informację na rzeczywistym
sprzęcie: redukcja błędu $\langle O\rangle$ jest najsilniejsza przy niskim koszcie CZ i słabnie
monotonicznie z routingiem, natomiast odzysk dokładności włącza się dopiero powyżej progu
szumu  --- a przy tym DD na tym urządzeniu szkodzi
mimo że symulator przewiduje korzyść, a niezależny od zadania model mirror fidelity wiąże
wszystko z błędami pojedynczych bramek.

% =====================================================================
\section{Uzasadnienie rezygnacji z kubitów logicznych (QEC)}
\label{sec:qec}
% =====================================================================

Rozważałam zakodowanie każdego kubitu logicznego w około 13 kubitów fizycznych (rząd
wielkości kodu powierzchniowego $d=3$). Odrzucam pełne QEC z trzech niezależnych
powodów~\cite{li2017,cai2023}:

\begin{enumerate}
\item \textbf{Budżet kubitów.} $5 \times 13 = 65 > 54$ (Emerald) --- nie mieści się nawet na
      największym urządzeniu z planu.
\item \textbf{Wymóg poniżej progu.} QEC pomaga tylko wtedy, gdy błąd fizyczny jest mniejszy
      niż próg (około 0,5--1\%). ODRA ma $p_{2q}\approx 1{,}9\%$, a Garnet i Emerald są blisko
      progu; w tym reżimie kodowanie dokłada więcej błędu, niż usuwa.
\item \textbf{Bramki nie-Clifforda.} VQC używa ciągłych obrotów zależnych od danych
      ($R_y(x)$), nietranswersalnych --- wymagają one kosztownej destylacji stanów
      magicznych~\cite{li2017}.
\end{enumerate}

„Wariant logiczny” realizuję zatem jako uzasadnioną wersję zredukowaną (Deliverable D3).

% =====================================================================
\section{Deliverables, podsumowanie i wnioski}
\label{sec:deliverables}
% =====================================================================

\subsection{Status deliverables Miesiąca 1}

\begin{longtable}{@{}p{5.4cm}cp{7.6cm}@{}}
\toprule
\textbf{Deliverable} & \textbf{Status} & \textbf{Gdzie w dokumencie / w kodzie} \\
\midrule
\endhead
D1. Przegląd literatury (mitygacja ${+}$ narzut kubitów logicznych)
  & \checkmark
  & Sekcja~\ref{sec:literatura} oraz \emph{Literatura}; kontekst teoretyczny w sekcji 2 \\
\addlinespace
D2. Raport baseline'u obwodu i metryk
  & \checkmark
  & Sekcje 5.1--5.6 (tabele i rysunki); kod: \texttt{freeze\_reference.py} (zamrożona
    referencja bezszumowa), \texttt{train\_vqc.py}, \texttt{train\_patterns.py};
    dane: \texttt{review\_notes.md}, katalog \texttt{figures/} \\
\addlinespace
D3. Wariant zakodowany/logiczny albo uzasadniona wersja zredukowana
  & \checkmark
  & Sekcja~\ref{sec:qec} (analiza QEC $\to$ QEM, trzy niezależne powody odrzucenia) \\
\addlinespace
D4. Protokół eksperymentalny mapowania optymalnego i niekorzystnego
  & \checkmark
  & Sekcja 5.1 (definicja routing ladder, tabela~\ref{tab:tc-ladder}); kod:
    \texttt{routing\_sweep.py} (generacja routing ladder), \texttt{run\_hardware.py}
    (wykonanie na QPU) \\
\bottomrule
\end{longtable}

\subsection{Podsumowanie i wnioski}

\begin{itemize}
\item Zbudowałam powtarzalny zestaw narzędzi QEM z treningiem w kalibrowanym symulatorze
      i inferencją na sprzęcie. Zbiór \texttt{two\_curves} został zmierzony na rzeczywistym
      QPU zarówno pod względem błędu $\langle O\rangle$, jak i mirror fidelity, dzięki czemu
      stanowi gotowy punkt odniesienia --- ten sam protokół można teraz zastosować bez zmian
      do ewaluacji na większych urządzeniach (Garnet, Emerald) i porównać wyniki wprost.
\item Błąd $\langle O\rangle$ rośnie monotonicznie z kosztem CZ, a odzysk ZNE maleje z CZ ---
      spójnie na symulatorze i na sprzęcie. Odpowiada to na pytania badawcze 3 i 4.
\item \textbf{Skuteczność ZNE maleje monotonicznie z kosztem CZ.} Na pełnym foldzie
      współczynnik redukcji błędu $\langle O\rangle$ spada z $3{,}7\times$ przy 16 CZ do
      $1{,}04\times$ przy 87 CZ (mitygacja praktycznie nieskuteczna). Bezpiecznym wnioskiem
      jest \emph{malejąca opłacalność mitygacji w okolicach $85$--$90$ CZ}, a nie aktywne
      pogarszanie dokładności. Ilościowe wyznaczenie tego progu na
      Garnecie i Emeraldzie (inna topologia, inny profil błędów) jest zadaniem Miesiąca 2.
\item \textbf{ZNE odzyskuje dokładność klasyfikatora na sprzęcie} tam, gdzie szum przerzucił
      znak $\langle O\rangle$: baseline przy 52 CZ zyskuje $0{,}706\to0{,}767$ (macro-F1
      równolegle, bez zapaści klasowej). Przy 87 CZ odzysk jest już w granicach rozrzutu
      międzyfoldowego (walidacja 5-fold: $0{,}626\to0{,}643$), co wyznacza górną granicę
      opłacalności mitygacji.
\item \textbf{DD w natywnej konfiguracji urządzenia jest wynikiem negatywnym} ---
      pogarsza wszystkie mierzone wielkości w czterech niezależnych eksperymentach.
      Właściwym opisem zastosowanego potoku jest zatem REM${}+{}$ZNE.
\item \textbf{Porównanie z pracą źródłową (raw wobec raw, walidacja 5-fold) nie odtwarza
      raportowanej zapaści pierścieni przy L=6}: native pozostaje zdrowy, baseline degraduje się
      bez zapaści klasowej, a różnicę --- ponieważ wagi, ansätze i transpilacja są identyczne ---
      przypisuję dryfowi kalibracji urządzenia. Wynik wzmacnia tezę pracy o zależności zachowania
      on-device od kalibracji.
\end{itemize}

% =====================================================================
\section{Repozytorium}
% =====================================================================

Kod projektu jest dostępny na GitHub:
\url{https://github.com/karinaleskiewicz06/QEM}

\medskip
\noindent Główne pliki:

\begin{itemize}
\item \texttt{ansatz.py} --- HEA z data re-uploading, wzorce entangling i liczenie
      kosztu dwukubitowego. Testy: \texttt{test\_ansatz.py}.
\item \texttt{noisy\_sim.py} --- kalibrowany model szumu Aer (gwiazda ODRA), harmonogram
      bezczynności, REM, ładowanie kalibracji.
\item \texttt{zne\_dryrun.py} --- prymitywy ZNE (unitary folding, ekstrapolacja
      liniowa / Richardsona / wykładnicza), przebieg ablacyjny w symulatorze.
\item \texttt{run\_hardware.py} --- wykonanie \texttt{two\_curves} na ODRA (iqm-client),
      staged rollout, wznawialny zapis, ablacja DD.
\item \texttt{collect\_odra\_calibration.py} --- pobranie i remap kalibracji ODRA.
\end{itemize}

\paragraph{Odtworzenie wyników (skrypt $\to$ dane $\to$ tabela/rysunek).} Każdy wynik w
raporcie da się prześledzić do konkretnej komendy i pliku:

Skrypty z prefiksem \texttt{banknote\_study/} podano bez tego prefiksu dla zwięzłości.

\begin{center}
\scriptsize
\renewcommand{\arraystretch}{1.25}
\begin{tabular}{@{}>{\raggedright\arraybackslash}p{5.1cm}>{\raggedright\arraybackslash}p{5.3cm}>{\raggedright\arraybackslash}p{3.3cm}@{}}
\toprule
\textbf{Skrypt / komenda} & \textbf{Plik wynikowy} & \textbf{W raporcie} \\
\midrule
\texttt{routing\_sweep.py} & \texttt{routing\_ladder.csv} & Tab.~\ref{tab:tc-ladder} (5.1) \\
\texttt{run\_hardware.py -\/-submit} & \texttt{hardware\_results.csv} & Sekcja 5.2 \\
\texttt{mirror\_fidelity\_two\_curves.py} & \texttt{two\_curves\_mirror\_} \texttt{fidelity\{,\_sim\}.csv} & Sekcja 5.3 \\
\texttt{noiseless\_and\_cost.py} & \texttt{banknote\_noiseless\_} \texttt{cost.csv} & Sekcja 5.4 \\
\texttt{run\_zne\_banknote.py} & \texttt{banknote\_zne\_} \texttt{d\{2,4,6\}.csv} & Sekcja 5.5 \\
\texttt{run\_hardware\_banknote.py} \texttt{-\/-submit} & \texttt{banknote\_hardware\_} \texttt{fullfold.csv} & Tab.~\ref{tab:bank-qpu-full} (5.6) \\
ten sam, \texttt{-\/-folds 1..5} \texttt{-\/-depths 6} & \texttt{banknote\_hardware\_} \texttt{cv\_d6.csv} & 5.6 (CV, w toku) \\
\texttt{mirror\_fidelity.py} & \texttt{banknote\_mirror\_} \texttt{fidelity.csv} & Sekcja 5.7 \\
\texttt{make\_figures.py} & \texttt{figures/*.png} & rysunki sekcji 5 \\
\bottomrule
\end{tabular}
\renewcommand{\arraystretch}{1.0}
\end{center}

\noindent Skrypty sprzętowe domyślnie działają w trybie \texttt{-\/-dry-run} (budują i walidują
obwody lokalnie, nic nie wysyłają); wysyłka wymaga flagi \texttt{-\/-submit} oraz zmiennych
\texttt{IQM\_SERVER\_URL} i \texttt{IQM\_TOKEN}. Zapis jest wznawialny (\texttt{-\/-resume}).

% =====================================================================
\section{Plan na Miesiąc 2}
% =====================================================================

Na podstawie wyników Miesiąca 1 plan na Miesiąc 2 obejmuje trzy kierunki:

\begin{itemize}
\item \textbf{Rozszerzenie na pozostałe dwa urządzenia.} Przeniesienie routing ladder na
      IQM Garnet (20 kubitów) i IQM Emerald (54 kubity), czyli na topologię kraty
      kwadratowej, aby sprawdzić, czy obserwowane zależności są własnością gwiazdy, czy
      ogólną własnością narzutu routingu.
\item \textbf{Wyznaczenie progu.} Ilościowe określenie progu kosztu dwukubitowego, przy
      którym mitygacja przestaje być opłacalna --- na podstawie zaobserwowanego w Miesiącu 1
      słabnięcia ZNE w okolicach $85$--$90$ CZ, zmierzonego teraz na trzech różnych
      urządzeniach i topologiach.
\item \textbf{Warstwa „logical-ish”.} Detekcja błędów z post-selekcją (małe kody, na przykład
      $[[4,2,2]]$) jako bliskoterminowy odpowiednik ochrony logicznej --- w miarę możliwości
      sprzętowych, bez wymogu szybkiego feed-forward.
\end{itemize}

% =====================================================================
\section{Przegląd literatury}
\label{sec:literatura}
% =====================================================================

Poniżej adnotowany przegląd, pogrupowany według filarów projektu: mitygacja błędów
(ZNE, REM), narzut kubitów logicznych / QEC oraz benchmarking QML (zbiór danych, encoding,
ansatz).

\paragraph{Mitygacja błędów --- przegląd.}
Cai i in., \emph{Quantum Error Mitigation}~\cite{cai2023}. Kompletny przegląd metod QEM
w erze NISQ (ZNE, PEC, REM, symetryzacja) wraz z ramami estymatora: dekompozycja MSE na bias
i wariancję oraz narzut próbkowania $C_{em}$. \textbf{Związek z projektem:} dostarcza
formalizmu, w którym mierzę błąd $\langle O\rangle$ (redukcja biasu) i uzasadniam, dlaczego
dokładność jest metryką wtórną; punkt odniesienia dla wyboru potoku mitygacji.

\paragraph{Zero-Noise Extrapolation.}
Li i Benjamin, \emph{Efficient variational quantum simulator incorporating active error
minimisation}~\cite{li2017}. Jedna z prac wprowadzających ideę sztucznego zwiększania błędu
i ekstrapolacji do zera; omawia też, dlaczego QEC jest kosztowne (zwielokrotnienie zasobów,
destylacja stanów magicznych). \textbf{Związek z projektem:} teoretyczne źródło mojego ZNE
oraz część uzasadnienia rezygnacji z pełnego QEC (sekcja~\ref{sec:qec}).

Giurgica-Tiron i in., \emph{Digital zero noise extrapolation}~\cite{giurgica2020}. Wprowadza
unitary folding ($U \to U(U^{\dagger}U)^k$) jako cyfrowe skalowanie szumu dostępne na poziomie
bramek oraz porównuje ekstrapolacje: liniową, Richardsona i wykładniczą. \textbf{Związek
z projektem:} bezpośrednia podstawa mojej implementacji ZNE (skale $\{1,3,5\}$, trzy warianty
ekstrapolacji), wykonanej samodzielnie w repozytorium.

\paragraph{Readout Error Mitigation.}
Koh, Koh i Thompson, \emph{Readout Error Mitigation for Mid-Circuit Measurements and
Feedforward}~\cite{koh2024}. Formalizuje błędy odczytu jako proces klasyczny i podaje protokół
REM (około 60\% redukcji błędu na sprzęcie nadprzewodzącym), także dla pomiarów
śródobwodowych. \textbf{Związek z projektem:} uzasadnia REM jako warstwę bazową przed ZNE;
stosuję dokładny, per-kubitowy model $Z/I$ ($\langle Z\rangle/(1-2p)$) zamiast pełnej
macierzy $A$.

\paragraph{Benchmarking QML --- zbiór danych, encoding, ansatz.}
Bowles, Ahmed i Schuld, \emph{Better than classical?}~\cite{bowles2024}. Duże, systematyczne
badanie 12 modeli QML na 6 zadaniach (pakiet \texttt{qml-benchmarks}, źródło
\texttt{two\_curves}); pokazuje między innymi, że usunięcie splątania często nie szkodzi.
\textbf{Związek z projektem:} uzasadnia wybór \texttt{two\_curves}, schemat angle encoding
z data re-uploading oraz traktowanie wzorca entangling jako obiektu badań, a nie źródła
przewagi.


% =====================================================================
\begin{thebibliography}{9}
% =====================================================================

\bibitem{cai2023}
Z. Cai, R. Babbush, S. C. Benjamin, S. Endo, W. J. Huggins, Y. Li, J. R. McClean,
T. E. O'Brien, \emph{Quantum Error Mitigation}, arXiv:2210.00921 (2023).

\bibitem{giurgica2020}
T. Giurgica-Tiron, Y. Hindy, R. LaRose, A. Mari, W. J. Zeng, \emph{Digital zero noise
extrapolation for quantum error mitigation}, IEEE QCE / arXiv:2005.10921 (2020).

\bibitem{koh2024}
J. M. Koh, D. E. Koh, J. Thompson, \emph{Readout Error Mitigation for Mid-Circuit
Measurements and Feedforward}, arXiv:2406.07611 (2024).

\bibitem{bowles2024}
J. Bowles, S. Ahmed, M. Schuld, \emph{Better than classical? The subtle art of benchmarking
quantum machine learning models}, arXiv:2403.07059 (2024).

\bibitem{li2017}
Y. Li, S. C. Benjamin, \emph{Efficient variational quantum simulator incorporating active
error minimisation}, Phys. Rev. X \textbf{7}, 021050 / arXiv:1611.09301 (2017).

\bibitem{ondevice2026}
\emph{On-Device Verification of Variational Quantum Classifiers}, zgłoszenie anonimowe
(2026). Źródło modeli i danych Banknote:
\url{https://github.com/Kolo-Naukowe-Axion/QC1_Quantum_Banknote_Classifier_on_ODRA_5}.

\end{thebibliography}

\end{document}
